% =====================================================================
%  main_short.tex  --  CONDENSED VERSION
%
%  Target: CISSE 30th Colloquium, deadline 15 September 2026.
%  CISSE requires 3,000-4,000 words including a <=250 word abstract.
%  The full-length main.tex runs ~5,500 body words and is over limit;
%  this file is the submission version.
%
%  IMPORTANT: CISSE requires the official IEEE conference template.
%  On Overleaf, comment out the \documentclass line below and use
%  \documentclass[conference]{IEEEtran} instead. The article-class
%  fallback here exists only so the file compiles without IEEEtran
%  installed; it approximates but does not equal the IEEE template.
%
%  Also required by the CFP:
%    - anonymized for double-blind review
%    - ORCID iDs for all authors in the final submission
%    - DOI links in references where available
%    - figures at minimum 300 DPI
%    - IRB / ethics approval information for human-subject research
%    - disclosure of generative AI or automated tool use
%    - student authors must declare student status at submission
% =====================================================================

\documentclass[10pt,twocolumn,letterpaper]{article}
% \documentclass[conference]{IEEEtran}   % <-- REQUIRED for CISSE

\usepackage[letterpaper,margin=0.72in,columnsep=0.25in]{geometry}
\usepackage{mathptmx}
\usepackage[T1]{fontenc}
\usepackage[utf8]{inputenc}
\usepackage{graphicx}
\usepackage{booktabs}
\usepackage{tabularx}
\usepackage{array}
\usepackage{xcolor}
\usepackage{enumitem}
\usepackage{caption}
\usepackage{balance}
\usepackage{titlesec}
\usepackage[hidelinks,breaklinks]{hyperref}
\usepackage{url}

\captionsetup{font=small,labelfont=bf,skip=3pt}
\setlist{nosep,leftmargin=1.1em}
\titleformat{\section}{\normalfont\large\bfseries}{\thesection.}{0.5em}{}
\titleformat{\subsection}{\normalfont\normalsize\bfseries}{\thesubsection.}{0.5em}{}
\titlespacing*{\section}{0pt}{1.1ex}{0.6ex}
\titlespacing*{\subsection}{0pt}{0.9ex}{0.4ex}
\renewcommand{\arraystretch}{1.1}
\newcommand{\ig}[1]{\textsc{ig#1}}
\newcommand{\todo}[1]{\textcolor{red}{[\textbf{TODO:} #1]}}

\title{\vspace{-1.6em}\bfseries From Threats to Controls: A CIS-Grounded
Solutions Dimension and Gameplay Framework for the Security Cards}
\author{\normalsize Anonymous Submission}
\date{}

\begin{document}
\maketitle

\begin{abstract}
\noindent\small
The Security Cards organize security threat brainstorming into four dimensions
and deliberately stop at threat discovery. In cybersecurity clinic courses,
where students advise under-resourced community organizations, discussion that
ends at threat identification leaves the central professional task---recommending
proportionate controls under real budget constraints---unpracticed. We present a
fifth dimension, \emph{Solutions}: twenty-one defensive control cards mapped to
the CIS Critical Security Controls v8.1 at the safeguard level, plus seven
zero-cost modifier cards carrying architectural and risk-treatment principles.
We follow a design-based research process across three cycles: a formative study
of the initial deck, a redesign in which every change traces to a requirement
that study produced, and a controlled classroom evaluation. Card cost is derived
from CIS Implementation Group tier rather than designer judgment, and the budget
is calibrated against expected hand value so that trade-offs are forced rather
than nominal. A coverage analysis against
Implementation Group 1 identified five foundational gaps in our initial
fifteen-card deck---asset inventory, software inventory, secure configuration,
email and browser protections, and incident response planning---which motivated
the current version and replaced designer intuition with an auditable procedure.
We contribute the mapped deck, a six-variant gameplay framework, and a core loop
in which a risk matrix is scored before and after control selection so that
residual risk is simultaneously a game mechanic and a measurable artifact.
Cycle 1, with 19 multidisciplinary participants, established feasibility and
accessibility and produced four design requirements; we claim no learning gains
from it. Cycle 3, a controlled classroom study, is scheduled for September 2026
and its design is reported in full. The deck is released under an open license.

\vspace{0.5em}
\noindent\textbf{Keywords:} cybersecurity education, game-based learning,
serious games, CIS Controls, design-based research, cybersecurity clinics
\end{abstract}

\vspace{0.6em}

% =====================================================================
\section{Introduction}

Cybersecurity education faces a structural problem: the material changes faster
than curricula can be rewritten, and the reasoning practitioners actually
perform---weighing partial information, arguing a position, committing resources
under constraint---is poorly served by lecture and multiple-choice assessment.
Active and collaborative methods address this gap in engineering education
generally \cite{prince2004active}, and game-based learning has become one
practical vehicle in security specifically \cite{plass2015gbl,wouters2013meta}.

Our setting is the cybersecurity clinic, a skill-based program in which
university students deliver security services to under-resourced organizations
in their community \cite{clinics_consortium}. Clinic students face an unusual
demand curve: within one semester they move from little or no security
background to advising a nonprofit director or municipal clerk on what to do
next, given a budget frequently near zero. The early weeks therefore need an
activity that requires no prerequisite vocabulary, exercises prioritization
rather than recall, and ends in a defensible recommendation rather than a list
of worries.

The Security Cards \cite{denning2013securitycards} meet the first two
requirements well. The 42-card toolkit organizes security thinking into four
dimensions---Human Impact, Adversary's Motivations, Adversary's Resources,
Adversary's Methods---and widens the space of threats a group considers. It is,
by its authors' framing, a brainstorming toolkit rather than a game: no win
condition, no scoring, no resource constraint. That is an appropriate choice for
its stated purpose, and it is the source of the third gap. In our use,
card-driven discussion reliably produced rich threat enumeration and then
stopped. Students left with a longer list of things that could go wrong and no
practiced vocabulary for what to do about it.

This paper contributes the missing half of that loop: (1) a Solutions dimension
of twenty-eight cards mapped to CIS Controls v8.1 \cite{cis2024controls} at
safeguard level, not merely control titles; (2) a cost model derived from CIS
Implementation Groups \cite{cis_ig} rather than designer intuition, with the
budget calibrated so the constraint binds; (3) a coverage analysis that
identified five foundational omissions in our own earlier deck and supplied a
non-arbitrary rationale for the current card set; (4) a six-variant gameplay
framework with a core loop that scores residual risk before and after control
selection.

\textbf{Method.} We report this as design-based research, in which an artifact
and the requirements it satisfies are developed together across iterative cycles.
\textbf{Cycle 1} (Section~\ref{sec:cycle1}) evaluated a fifteen-card deck with 19
participants; its purpose was to generate design requirements, not to validate a
finished artifact. \textbf{Cycle 2} (Section~\ref{sec:cycle2}) is the redesign,
in which each change traces to a requirement Cycle 1 produced. \textbf{Cycle 3}
(Section~\ref{sec:planned}) is a controlled classroom evaluation scheduled for
September 2026, reported here as a design so that it can be criticized before
rather than after data collection.

We are explicit about what we do not claim. Cycle 1 establishes that the artifact
is usable by novices in under ten minutes and identifies what was wrong with it.
It does not establish learning gains, and the Cycle 2 deck has not yet been
played in a classroom.

% =====================================================================
\section{Background and Related Work}

Two ideas are frequently conflated. \emph{Gamification} applies game design
elements to contexts that are not games \cite{deterding2011gamification};
\emph{game-based learning} uses an actual game as the learning environment
\cite{plass2015gbl,landers2014gamified}. This is game-based learning work: we
add rules, roles, scarcity, and scoring to a toolkit that previously had none.
Beyond self-reported engagement, game-based learning environments have been shown
to produce measurable differences in frontal brain activity relative to
conventional instruction \cite{kober2020frontal}.

Security education has a developed tradition of card and tabletop games
\cite{shostack_games}. \textbf{Control-Alt-Hack}
\cite{denning2013controlalthack} targets novices with a dispositional aim,
conveying that security work exists and is interesting rather than teaching an
analytic procedure. \textbf{Elevation of Privilege} \cite{shostack2014eop}
converts STRIDE threat modeling into a card game but assumes a system diagram
and a developer audience, placing it downstream of where clinic students begin.
\textbf{[d0x3d!]} \cite{gondree2013d0x3d} builds intuition for
attacker--defender dynamics on a board. \textbf{Riskio} \cite{hart2020riskio}
addresses non-technical organizational staff but requires an experienced Game
Master. \textbf{Backdoors \& Breaches} \cite{young2021backdoors,bhis2019backdoors} is the closest
analogue to our Solutions dimension---players select procedure cards against a
hidden attack---but its focus is incident response rather than \emph{ex ante}
control selection, and its cards are detection techniques rather than a
framework-aligned control catalogue. \textbf{MalAware} \cite{malaware2024}
structures a malware tabletop around NIST CSF functions, again with a response
emphasis. \textbf{Cyber Defence Dice} \cite{cyberdefencedice2026} pursues a
complementary goal, minimizing setup and play time for entry-level awareness,
and its survey identifies facilitator dependence and session length as recurring
barriers.

The gap is specific. Existing games teach learners to \emph{find} threats or to
\emph{respond} to incidents. None asks a novice to select a proportionate set of
preventive controls against a stated budget and justify the choice against a
recognized framework. That is the task a clinic student performs for a client in
week eight.

The CIS Critical Security Controls v8.1 organize defensive practice into 18
controls containing 153 safeguards, partitioned into three Implementation
Groups \cite{cis2024controls,cis_ig}. \ig{1} defines essential cyber hygiene:
the 56 safeguards a small organization with limited expertise should implement
first. Three properties make this an unusually good fit. \ig{1} is defined
\emph{for} the class of organization clinic students serve. The tiering supplies
an externally justified cost ordering, removing the need for designer-assigned
point values. And control identifiers give students a durable reference they can
pursue after the game ends.

% =====================================================================
\section{Cycle 1: The Initial Deck and Its Formative Study}
\label{sec:cycle1}

Cycle 1 evaluated a fifteen-card deck using a single variant: teams were given a
scenario and asked to rank the three most important defensive measures, with no
budget and no risk matrix. Its purpose was to determine whether novices can use
the deck without facilitator expertise, how long play takes, and whether the
vocabulary is accessible. We claim no learning gains; a single session with a
post-play survey and no control condition cannot support such a claim.
\todo{Insert IRB protocol number or exemption determination.}

Nineteen participants were recruited from a summer research program at an
undergraduate institution, spanning psychology, biology, physics, chemistry,
mathematical sciences, and computer science; Table~\ref{tab:participants} reports
the composition.
\todo{Complete from your records: reviewers specifically requested discipline
distribution, level, and prior gaming experience.}
Six teams of two to three played simultaneously, each facilitated by a member of
the research team. Teams received a hypothetical organization and an attack
scenario, then used the deck to select and rank the three most important
defensive measures. Sessions ran seven to ten minutes excluding debrief and
survey. Play was collaborative, and the budget mechanic was not yet implemented.

\begin{table}[h]
\centering
\caption{Formative evaluation participants ($N=19$).}
\label{tab:participants}
\footnotesize
\begin{tabular}{@{}l r l r@{}}
\toprule
\textbf{Discipline} & \textbf{$n$} & \textbf{Level / background} & \textbf{$n$} \\
\midrule
Computer science & \textcolor{red}{--} & Undergraduate & \textcolor{red}{--} \\
Psychology & \textcolor{red}{--} & Graduate & \textcolor{red}{--} \\
Biology & \textcolor{red}{--} & Faculty & \textcolor{red}{--} \\
Physics & \textcolor{red}{--} & No prior security course & \textcolor{red}{--} \\
Mathematical sciences & \textcolor{red}{--} & One security course & \textcolor{red}{--} \\
Other & \textcolor{red}{--} & More than one & \textcolor{red}{--} \\
\midrule
\multicolumn{4}{@{}l}{Self-rated security knowledge (1--5): mean 2.53, SD \textcolor{red}{--}}\\
\bottomrule
\end{tabular}
\end{table}

Seventeen of nineteen agreed the cards helped them learn. Ease of use was the
highest-rated item at 4.42 on a five-point scale; perceived sufficiency of card
content scored 3.95.
Figure~\ref{fig:likert} reports the full response distributions.

\begin{figure}[h]
\centering
\includegraphics[width=\columnwidth]{likert.pdf}
\caption{Post-play survey responses, Cycle 1 ($N=19$), reported as counts.}
\label{fig:likert}
\end{figure}

Open-ended responses converged on three themes. Participants valued the
definitions and examples on the card backs and the non-specialist vocabulary,
speaking directly to G1. Participants wanted feedback on whether their rankings
were correct, and noted that several solutions seemed equally valid.
Participants found fifteen cards to be many, with some descriptions overlapping.

\subsection{Design Requirements Produced by Cycle 1}
\label{sec:reqs}

Cycle 1 establishes two positive findings and yields four requirements. The
positive findings are that the activity is feasible in under ten minutes, which
matters in crowded syllabi, and that it is usable by participants with a mean
self-rated knowledge of 2.53 and no facilitator expertise beyond reading the
scenario aloud. It establishes nothing about learning.

The requirements follow directly from the reported shortcomings:

\begin{description}[leftmargin=1.6em,style=nextline]
\item[R1. Choices must be contestable.] Participants reported that many
responses seemed equally acceptable and that there was no clear distinction
between stronger and weaker choices. Unconstrained ranking has no wrong answers,
so it generates agreement rather than argument---and disagreement is the
mechanism by which the activity teaches.
\item[R2. Bound the working set.] Participants found fifteen cards to be many
and some descriptions overlapping. Facilitators independently found that
evaluating one dimension at a time reduced overload.
\item[R3. Provide a basis for feedback on quality.] Participants wanted guidance
on whether their rankings were appropriate. A defensible answer key requires
that reasoning be recorded in a scoreable form.
\item[R4. Make the card set auditable.] This requirement came from peer review
rather than from participants: asked why these fifteen cards, the honest answer
was designer intuition informed by clinic experience.
\end{description}

Section~\ref{sec:cycle2} addresses each in turn. We note that R1 and R3 are the
same observation seen from two sides, and that a single mechanism---making
choices cost something---answers both.

% =====================================================================
\section{Cycle 2: The Revised Deck}
\label{sec:cycle2}
\label{sec:design}

Every change in this cycle traces to a requirement from
Section~\ref{sec:reqs}: the budget and cost model answer \textbf{R1}, the
tiering and eight-card hand answer \textbf{R2}, the justification rule and
scoreable artifacts answer \textbf{R3}, and the coverage analysis answers
\textbf{R4}. Four standing design goals also carry over, the first two inherited
from the original toolkit: \textbf{G1} vocabulary accessibility, so a participant with no prior
coursework can act on a card without explanation; \textbf{G2} discussion
generation, since disagreement is the mechanism by which the activity teaches;
\textbf{G3} framework traceability, so every claim made at the table is
traceable to a published safeguard; and \textbf{G4} constraint realism, because
the defining feature of the organizations clinic students serve is that they
cannot do everything.

We preserved the double-sided format and dimension coloring of the original
deck, replaced its evocative photographs with pictographic icons after
facilitators observed participants interpreting the photographs literally, added
an \emph{Artificial Intelligence} card to Adversary's Resources, and moved the
Solutions palette away from the existing Human Impact blue, adding a distinct
icon frame so the dimensions separate in grayscale and for color-vision-deficient
players. We report the icon change as a design decision motivated by facilitator
observation, not as an empirical finding; it is confounded with other
simultaneous revisions, and Section~\ref{sec:planned} includes a comprehension
check designed to test it properly.

Each Solutions card names a defensive practice, states what it does in one
sentence of non-specialist language, gives two concrete examples, cites the CIS
Control and representative safeguards, and carries a discussion prompt.
Table~\ref{tab:mapping} presents the mapping. Two properties matter. Mapping is
at safeguard level: a card reading only ``CIS Control 6: Access Control
Management'' tells a novice nothing, whereas naming 6.1 and 6.2---grant access
on a documented process, revoke it when someone leaves---tells them what the
control asks an organization to do. And the \ig{} column and the cost column are
the same information; cost is not an aesthetic judgment but the Implementation
Group at which the card's representative safeguards sit.

\begin{table*}[t]
\centering
\caption{Solutions dimension: cards, CIS Control v8.1 mapping, representative
safeguards, Implementation Group, and derived play cost. Cards 9--13 were added
in the current version following the \ig{1} coverage analysis. Full card text is
available in the archived deck.}
\label{tab:mapping}
\footnotesize
\begin{tabularx}{\textwidth}{@{}r l X c c@{}}
\toprule
\textbf{\#} & \textbf{Card} & \textbf{CIS Control (representative safeguards)} &
\textbf{IG} & \textbf{Cost} \\
\midrule
\multicolumn{5}{@{}l}{\textit{Foundational tier}}\\
1  & Software Updates & 7 Continuous Vulnerability Mgmt (7.1, 7.3, 7.4) & 1 & 1 \\
2  & Social Engineering Awareness & 14 Security Awareness \& Skills Training (14.1, 14.2, 14.3) & 1 & 1 \\
3  & Deepfake \& AI Fraud Awareness & 14 Security Awareness \& Skills Training (14.2, 14.9) & 1--2 & 1 \\
4  & Identity \& Access Management & 6 Access Control Management (6.1, 6.2, 6.7) & 1--2 & 1 \\
5  & Password \& Authentication & 5 Account Mgmt (5.2, 5.4); 6 Access Control Mgmt (6.3--6.5) & 1 & 1 \\
6  & Backup \& Recovery & 11 Data Recovery (11.1--11.4) & 1 & 1 \\
7  & Account Lifecycle Management & 5 Account Management (5.1, 5.3, 5.5) & 1--2 & 1 \\
8  & Physical Security & No dedicated control; enables Controls 1 and 3 & --- & 1 \\
9  & Asset Inventory & 1 Inventory \& Control of Enterprise Assets (1.1, 1.2) & 1 & 1 \\
10 & Software Inventory & 2 Inventory \& Control of Software Assets (2.1, 2.2, 2.3) & 1 & 1 \\
11 & Secure Configuration & 4 Secure Configuration (4.1, 4.2, 4.6, 4.7) & 1 & 1 \\
12 & Email \& Browser Protections & 9 Email \& Web Browser Protections (9.1, 9.2, 9.6) & 1--2 & 1 \\
13 & Incident Response Planning & 17 Incident Response Management (17.1--17.4) & 1--2 & 1 \\
\midrule
\multicolumn{5}{@{}l}{\textit{Extended tier}}\\
14 & Encryption & 3 Data Protection (3.6, 3.10, 3.11) & 1--2 & 2 \\
15 & Audit Logging \& SIEM & 8 Audit Log Mgmt (8.2, 8.9, 8.11); 13 Network Monitoring (13.1) & 1--2 & 2 \\
16 & Endpoint Detection \& Response & 10 Malware Defenses (10.1, 10.6, 10.7) & 1--2 & 2 \\
17 & Third-Party \& Contractual Controls & 15 Service Provider Management (15.1, 15.4, 15.5) & 1--3 & 2 \\
18 & Network Segmentation & 12 Network Infrastructure Mgmt (12.2); 13 (13.4) & 2 & 2 \\
\midrule
\multicolumn{5}{@{}l}{\textit{Advanced tier}}\\
19 & Penetration Testing & 18 Penetration Testing (18.1, 18.2, 18.3, 18.5) & 2--3 & 3 \\
20 & Managed Detection \& Response & 8 (8.11); 13 (13.1, 13.11); 17 (17.1) & 2--3 & 3 \\
21 & Cyber Insurance \& Risk Transfer & No control; transfer is a risk treatment & --- & 3 \\
\midrule
\multicolumn{5}{@{}l}{\textit{Modifier cards (three dealt per round, at most two played; each requires a stated justification)}}\\
22 & Defense in Depth & Requires two purchases that fail by different means & --- & 0 \\
23 & Fail Safe Defaults & Requires naming the state a purchased control fails into \cite{saltzer1975protection} & --- & 0 \\
24 & Least Privilege & Requires card 4 or 7; name one access to reduce & --- & 0 \\
25 & Human Factors & Requires naming a purchase staff will circumvent & --- & 0 \\
26 & Risk Acceptance & Played instead of a purchase; requires a named decision owner & --- & 0 \\
27 & Compensating Control & Requires naming the unaffordable control and the residual gap & --- & 0 \\
28 & Deferred Investment & Caps spend at 4, banks 2 for a 10-point budget next engagement & --- & 0 \\
\bottomrule
\end{tabularx}
\end{table*}

\subsection{Coverage Analysis}
\label{sec:gaps}

\textbf{Addressing R4.} To replace intuition with a procedure, we asked, for
each of the 56 \ig{1} safeguards, whether any card in the deck would prompt a
player to consider it.

Our initial fifteen-card deck omitted five foundational areas while including
one outside \ig{1} entirely. Controls 1 and 2 (asset and software inventory)
were absent, though an organization cannot patch software it does not know it
runs. Control 4 (secure configuration) was absent, despite default credentials
being among the most common clinic findings. Control 9 (email and browser
protections) was absent even though phishing was the central threat in both of
our recorded sessions---students reached for awareness training as the only
anti-phishing lever because it was the only one present. Control 17 (incident
response) was absent, despite clinic students being who a client calls during an
incident. Meanwhile Control 18 (penetration testing) was present and is
\ig{2}/\ig{3} only, inconsistent with our stated \ig{1} orientation; we retained
it, because clients ask about it, but re-tiered it to cost~3 so its expense is
visible in play.

Cards 9--13 close these gaps, all at \ig{1} and cost~1, shifting the deck's
center of gravity toward essential cyber hygiene as originally intended but not
originally achieved. We deliberately added nothing for Control 16 (application
software security), which begins at \ig{2} and addresses in-house development
that these clients rarely undertake. We report the analysis because the
procedure, not the resulting card list, is the transferable contribution: any
deck extension of this kind can be audited against a published framework rather
than defended on the designer's taste.

\subsection{Calibrating the Budget}
\label{sec:budget}

\textbf{Addressing R1.} A budget makes choices contestable, but only if it is
small relative to what is on the table. Our first implementation dealt eight
cards against ten points, which on inspection
is barely a constraint. The twenty-one purchasable cards carry 32 points, a mean
of 1.52 per card, so a random eight-card hand is worth roughly 12.2 points. Ten
points buys 82\% of such a hand, and a hand drawn entirely from the Foundational
tier is fully purchasable---a team in that position makes no trade-off at all.
We therefore set the default at \textbf{six points}, which buys roughly half a
typical hand and makes any Advanced-tier card cost half of everything a team has
(Table~\ref{tab:budget}).

\begin{table}[h]
\centering
\caption{Budget calibration against an eight-card hand worth $\approx$12.2
points. Larger budgets model better-resourced organizations.}
\label{tab:budget}
\footnotesize
\begin{tabular}{@{}c c c l@{}}
\toprule
\textbf{Budget} & \textbf{\% hand} & \textbf{Cost-3 card} & \textbf{Organization modeled} \\
\midrule
5  & 41\% & 60\% & Volunteer-run, no IT staff \\
\textbf{6}  & \textbf{49\%} & \textbf{50\%} & \textbf{Default; \ig{1}} \\
10 & 82\% & 30\% & \ig{2}; one IT staff member \\
15 & --- & 20\% & \ig{3}; security function \\
\bottomrule
\end{tabular}
\end{table}

Budget is therefore a property of the \emph{organization} in the scenario rather
than of the game, and tying it to Implementation Group makes that explicit: the
same threat facing a volunteer-run food bank and a county health department
warrants different answers, and students should feel that difference rather than
be told about it. For the study in Section~\ref{sec:planned} we hold the budget
fixed at six, since varying it would confound the condition comparison.

\subsection{Modifier Cards}

\textbf{Addressing R3.} The seven modifier cards carry principles rather than
controls, cost nothing, and are each governed by one rule: a modifier may be played only if the team can
articulate the justification the card demands, or it returns to the deck. This
is what keeps a zero-cost card from being a free point, and it serves the
evaluation directly, since every modifier generates a recorded rationale.

The class exists because the original deck, and our first Solutions version,
taught exactly one risk treatment: buy a control. Practice recognizes
four---avoid, mitigate, transfer, accept---and a student who leaves believing
every risk must be fixed will give clients advice they cannot afford to take.
\emph{Cyber Insurance} supplies transfer, \emph{Risk Acceptance} supplies
acceptance, and \emph{Compensating Control} supplies the partial mitigation
under-resourced organizations actually implement. \emph{Deferred Investment}
adds a planning horizon: a team may cap spend at four of six and bank two,
receiving ten points---the \ig{2} budget---at the next engagement, while every
threat left unmitigated drifts one cell toward \emph{Very Likely}. Deferring is
how an organization climbs an Implementation Group, and the accrued risk is what
the climb costs.

Because the matrix gives likelihood only three states, a deferred threat reaches
\emph{Very Likely} within two engagements; rather than cap it there, a threat
surviving another engagement unmitigated \emph{realizes}. The outcome follows
not from chance but from what the team purchased earlier: discovery in days or
months depends on logging, data return on backups, knowing who to call on
incident response planning. This corrects a structural bias---twenty-one of
twenty-eight cards are preventive, and a team that never experiences an incident
never learns why detection and recovery were worth points.

Growth from fifteen to twenty-eight cards runs against a pilot finding that
fifteen already felt like many. Two mechanisms address this: no variant deals
more than eight control cards and three modifiers at once, and the tiers form a
curriculum sequence, with the thirteen Foundational cards alone supporting the
first weeks of a course.

% =====================================================================
\section{Gameplay Framework}
\label{sec:gameplay}

The original toolkit describes activities but not rules. We specify six variants
with explicit objectives, durations, and produced artifacts
(Table~\ref{tab:variants}), ordered by cognitive demand and designed to be run
in sequence across a semester or selected individually.

\begin{table}[h]
\centering
\caption{Gameplay variants ordered by cognitive demand. ``Artifact'' names the
observable product of play, which is what makes each variant assessable.}
\label{tab:variants}
\footnotesize
\begin{tabularx}{\columnwidth}{@{}l X l@{}}
\toprule
\textbf{Variant} & \textbf{Objective / artifact} & \textbf{Time} \\
\midrule
V1 Threat Model Generation & Build an incident narrative from four drawn threat elements; written narrative & 8--10 m \\
V2 Risk Matrix Placement & Estimate likelihood and severity; matrix coordinates with rationale & 10--15 m \\
V3 Solution Ranking & Prioritize among defensible alternatives; ranked top three & 7--10 m \\
V4 Control Mapping & Attribute an incident to the control that failed; attribution & 10 m \\
V5 Budgeted Defense & Allocate scarce resources; purchased portfolio and rationale & 15--20 m \\
V6 Residual Risk Loop & Recognize that controls reduce rather than remove risk; pre/post matrix delta & 20--25 m \\
\bottomrule
\end{tabularx}
\end{table}

V6 is the flagship and the variant we recommend for the classroom study. In
Phase~1 each team draws one card from each original dimension, constructs a
scenario, places the resulting threat on a $3\times3$ risk matrix (likelihood
$\times$ severity), and records the coordinates; severity placement must be
justified by naming the \emph{Human Impact} card in play, since novices reliably
rate severity by how sophisticated an attack sounds rather than by who is
harmed. In Phase~2 the team receives
eight Solutions cards, three modifiers, and six points, purchases a portfolio,
and states for each purchase which safeguard addresses which element of the
threat. In Phase~3 it re-places the same threat given the controls purchased and
states what would still get through.

This does three things at once. Pedagogically, the second placement is where
students discover that controls shift risk rather than remove it, and where the
team that spent everything on prevention must explain how it would detect what
it did not prevent. Mechanically, it closes the loop between risk assessment and
control selection. Methodologically, the pre/post delta and its accompanying
justifications are recorded artifacts scoreable against an expert key, which
converts a discussion activity into a measurable one without adding a test.

Teams of two to three play cooperatively against the scenario rather than
against each other, following the ICAP framework \cite{chi2014icap}, which
associates interactive engagement with deeper outcomes than merely active
engagement. The budget supplies the friction a cooperative structure otherwise
lacks: six points against eight cards makes it untrue that any answer will do.

Four optional facilitator-played \emph{inject} cards introduce adversity after
purchase---a grant not renewed, a vendor discontinuing a product, a key staff
member departing---each requiring a response rather than a loss the team can only
absorb. They are excluded from Cycle 3, since they introduce variance not held
constant across teams.

% =====================================================================
\section{Facilitation Findings}
\label{sec:facilitation}

Cycle 1 produced practical findings about running the activity that are
independent of the deck revision, and we report them because they transfer to
any instructor adopting either version.

\textbf{Pre-populate the risk matrix.} Placing two or three threats on the grid
before asking participants to position the rest substantially lowers the entry
cost of the matrix activity. A blank grid asks novices to establish a scale and
apply it simultaneously; a partly filled grid asks only the second.

\textbf{Evaluate one dimension at a time.} In ranking activities, considering
each dimension independently rather than all cards at once reduced apparent
cognitive load and produced more coherent discussion. This finding, from Cycle
1, is what the tiering and bounded hand size in Section~\ref{sec:cycle2}
generalize.

\textbf{Physical cards outperform digital, but digital works.} Preliminary
trials on a shared Zoom whiteboard indicate the activities survive the
transition to remote instruction, though physical cards generated visibly more
discussion. For programs teaching remotely, the activity remains viable; for
those with a choice, the physical deck is preferable.

\textbf{Facilitator expertise is not required.} Cycle 1 sessions were run by
researchers who read a scenario aloud and kept time. Cycle 2 adjudication is more
demanding, so the accompanying facilitator materials supply worked
accept-and-reject examples rather than relying on facilitator judgment.

\section{Cycle 3: Planned Classroom Evaluation}
\label{sec:planned}

Cycle 3 tests the Cycle 2 deck. We report its design in advance of running it,
so that it can be criticized before rather than after data collection and so
that the claims in this paper remain separable from the claims that study would
license. The
setting is an undergraduate cybersecurity clinic course during the first three
weeks of the Fall 2026 term.
\todo{Expected enrollment; IRB submission must precede the term.}

The design is a counterbalanced within-subjects crossover. Each team plays two
scenarios of matched difficulty: in one, teams use the four-dimension deck and
discuss defenses without Solutions cards; in the other, the five-dimension deck
with the budgeted V6 loop. Order and scenario--condition pairing are
counterbalanced, so each team serves as its own control and scenario effects are
separated from condition effects.

Measures are: \textbf{M1}, control-selection appropriateness, with purchased
portfolios scored against an expert key derived from CIS safeguard--threat
mappings by two raters blind to condition, agreement reported as Cohen's
$\kappa$ \cite{landis1977kappa}; \textbf{M2}, risk-matrix delta, the magnitude
and direction of movement between placements and whether stated residual risk is
consistent with controls purchased; \textbf{M3}, justification quality, scored
0--3 from bare assertion to a threat-specific mechanism correctly linked to a
safeguard; \textbf{M4}, mitigation breadth, the count of distinct measures
raised per scenario, which directly tests whether Solutions cards expand or
merely redirect discussion; \textbf{M5}, pre/post security self-efficacy
\cite{bandura1997selfefficacy}; \textbf{M6}, cognitive load via short-form
NASA-TLX \cite{hart1988nasatlx}, testing whether the expanded deck overloads
novices; and \textbf{M7}, a card-face-to-definition matching task testing the
icon revision.
\todo{Select a validated self-efficacy instrument rather than authoring items.}

We will use Wilcoxon signed-rank tests on M1--M4 given the expected sample size
and ordinal measures, with descriptive reporting for M5--M7, and will report
effect sizes and confidence intervals rather than significance alone. Four
threats to validity are anticipated and will be reported rather than controlled
for, since a single-course study cannot eliminate them: team size varies with
attendance, facilitators are not blind to condition, scenario difficulty is
matched by expert judgment rather than empirically, and novelty effects favor
whichever condition is played second.

% =====================================================================
\section{Discussion and Limitations}

The deck is an on-ramp, not a simulation. It does not carry the technical depth
required for a real risk assessment, and a student who plays it well is not
thereby competent to advise a client. Its function is to give a novice a
vocabulary, a framework identifier they can look up, and practice at the
specific cognitive move---from stated threat to proportionate, affordable
control---that the rest of a clinic curriculum builds on.

Anchoring to CIS Controls buys traceability and an externally justified cost
model, and costs generality: a course oriented to the NIST Cybersecurity
Framework \cite{nist_csf2} or working toward CSEC2017 knowledge areas
\cite{csec2017} would want a different mapping on the card backs. We regard the
mapping procedure in Section~\ref{sec:gaps} as portable even where the control
identifiers are not. A sixth Impact dimension was considered and
rejected: the original deck already contains one, and what was missing was not a
category but the moment at which consequence follows from the team's own
choices.

Cycle 1 had no control condition, no pre/post measure, and a single short
session, and it evaluated the fifteen-card deck rather than the deck described
in Section~\ref{sec:cycle2}. Under a design-based research framing this is the
expected role of a first cycle---to generate requirements---but it does mean no
empirical claim in this paper attaches to the revised deck. Participants were volunteers in a research program, which
likely selects for engagement. The icon revision is supported by facilitator
observation only. The cost model, though derived from Implementation Groups,
still involves judgment where a card spans tiers. And the current deck has not
yet been played in a classroom; Section~\ref{sec:planned} is a plan, not a
result.

Three directions seem worthwhile beyond Cycle 3. A sixth dimension representing
the \emph{defender's} resources would let scarcity vary by organization rather
than by rule. A scenario deck of named incident types would standardize scenarios
across sessions and remove a confound from any multi-team study. And
\emph{standing risks}---two chronic problems the client already has, placed on
the matrix alongside the presenting incident and competing for the same
budget---would test something the current single-threat loop cannot: whether
students recognize that a control touching several problems is worth more than
one addressing only today's. That is the argument for foundational controls, and
at present nothing in the design rewards it. All three are untested.

% =====================================================================
\section{Conclusion}

The Security Cards are effective for the task they were designed for and stop
deliberately at its boundary. In cybersecurity clinic education, where students
must convert threat understanding into affordable recommendations for
organizations that cannot afford much, that boundary falls in the wrong place.
We have presented a Solutions dimension of twenty-eight cards mapped to CIS
Controls v8.1 at safeguard level, a cost model derived from Implementation Group
tier and calibrated so the constraint binds, a coverage analysis that identified
and closed five foundational gaps in our own earlier deck, and a gameplay
framework whose core loop scores residual risk before and after control
selection. A formative evaluation established feasibility and accessibility; a
controlled study is scheduled for September 2026 and its design is reported
here. The deck is released openly, and we welcome its use, criticism, and
re-tiering by other clinic programs.
\todo{Replace with the Zenodo DOI; use an anonymous view-only link for the blind
submission.}

\section*{Generative AI Disclosure}
\todo{Required by the CFP. State what tools were used and for what. AI tools may
not be listed as authors.}

\input{references.tex}

\end{document}
